\documentclass[10pt,a4paper]{amsart}
\usepackage[margin=19mm]{geometry}
\usepackage{amsmath,amssymb,mathtools}
\usepackage[scr=rsfs,cal=euler]{mathalfa}
\usepackage{xfrac}
\usepackage[unicode,hidelinks,bookmarks=false]{hyperref}
\newcommand{\ie}{i.e.\ }
\newcommand{\cf}{cf.\ }
\newcommand{\resp}{respectively\ }
\newcommand{\Subsection}{\subsection}
\providecommand{\nobreakdash}{\nobreak\hbox{-}}
\setlength{\emergencystretch}{2em}
\multlinegap0pt
\usepackage{bm}
\usepackage{bbm}
\usepackage{dsfont}
\usepackage{xspace}
\usepackage{cancel}
\usepackage{mathtools}
\usepackage{amsthm}
\makeatletter
\@ifundefined{theorem}{\newtheorem{theorem}{Theorem}}{}
\@ifundefined{definition}{\newtheorem{definition}[theorem]{Definition}}{}
\@ifundefined{lemma}{\newtheorem{lemma}[theorem]{Lemma}}{}
\@ifundefined{remark}{\newtheorem{remark}[theorem]{Remark}}{}
\makeatother
\newcommand{\C}{\mathbb{C}}
\newcommand{\D}{\mathbb{D}}
\title[The Julia-Wolff-Carath\'eodory theorem in convex finite type]{The Julia-Wolff-Carath\'eodory theorem in convex finite type domains}
\author{Leandro Arosio, Matteo Fiacchi}
\date{Source: arXiv:2407.09199v2 (CC BY 4.0)}
\begin{document}
\maketitle

\noindent\emph{Source and licence.} The Julia-Wolff-Carath\'eodory theorem in convex finite type domains. Version arXiv:2407.09199v2 (CC BY 4.0), distributed under the Creative Commons Attribution 4.0 International licence, as recorded in the arXiv licence metadata for this exact version and shown on the abstract page https://arxiv.org/abs/2407.09199v2. Full rights notice: licenses/arxiv-2407.09199-CC-BY-4.0.txt.\par\smallskip

\noindent\textit{Editorial scope.} Counted unit: the Remark whose source label is \texttt{None} in the article (source0.tex lines 1302--1305, source label \texttt{None}), delivered together with the article's own complete original proof of it (source0.tex lines 1307--1328, one \texttt{proof} environment of 22 lines). No mathematics is written, repaired or re-derived: statement and proof are the article's own text line for line. Required context retained uncounted: restricted (lines 990--998); stronglyrestricted (lines 1236--1243); localLindelof (lines 1265--1267); goodprodev (lines 1285--1291). Every definition, display and label of those lines is retained unchanged. Omitted from this package and not redistributed: 1--989, 999--1235, 1244--1264, 1268--1284, 1292--1301, 1306--1306, 1329--2374. Nothing inside a retained excerpt is omitted or altered: every retained body is delivered exactly as the article's lines are written, so no omission receipt of any kind accompanies this package. Citations: the retained excerpts carry no citation command at all, so no bibliography entry of the article is transcribed and no reference is redistributed. Reference adaptation: none was needed and none was made; every reference inside the delivered unit resolves inside it, so no unresolved reference or citation is produced. Printed slips kept literal: no printed word, symbol, hypothesis or number of the counted statement or of its proof was altered. Layout and class: the article's own class is \texttt{amsart} with packages including amsfonts, amssymb, amscd, amstext, mathrsfs, color, xcolor, graphicx, tikz, hyperref, mathtools; the wrapper is the lane's pinned \texttt{amsart} editorial wrapper at 10pt on A4, which restates the theorem-like environments the retained text needs and adds the article's own macro definitions that the retained text needs (\texttt{\textbackslash C}, \texttt{\textbackslash D}), with extra wrapper packages (bm, bbm, dsfont, xspace, cancel, mathtools). The delivered numbering is the unit's own. Assets: no retained span contains a figure environment, a graphics-inclusion command, a PDF-page inclusion, a table or a TikZ picture; no image file is copied into this package, the package declares no asset dependency and no image file is redistributed with it. Licence: version arXiv:2407.09199v2 of this article is distributed under the Creative Commons Attribution 4.0 International licence, as recorded in the arXiv licence metadata for this exact version and shown on the abstract page https://arxiv.org/abs/2407.09199v2. A full rights notice is delivered at licenses/arxiv-2407.09199-CC-BY-4.0.txt. Characters: every retained excerpt is pure ASCII, so the delivered engine, XeLaTeX with its default Latin Modern text font, reports no missing character.
		\begin{definition}\label{restricted}
			Let $D\subset\C^d$ be a $\C$-proper convex domain and let $\xi\in\partial D$ be a point of locally finite type. Let $\varphi\colon\D\to D$ be a complex geodesic  with endpoint $\xi$, and let  $(z_n)$ be a sequence in $D$ converging to $\xi$.
			We say that  $(z_n)$  {\sl $K'$-converges} to $\xi$ (or is  {\sl restricted}) if it $K$-converges to $\xi$ and if
			\begin{equation}\label{eqrestricted}
				\lim_{n\to+\infty}k_D(z_n,\varphi(\D))=0.
			\end{equation}
			If $f\colon D\to\C$ is a function, we say that  $f$ has  {\sl $K'$-limit} (or {\sl restricted $K$-limit})  $L\in\C$ at $\xi$  if $f(z_n)\to L$ for any restricted sequence $(z_n)$ converging to $\xi$. In this case we write
			$$K'\textrm{-}\lim_{z\to\xi}f(z)=L.$$
		\end{definition}

		\begin{definition}
			A {\sl restricted curve} $\gamma:[0,1)\to D$ with endpoint $\xi$ is a curve such that for all sequences $(t_n)$ in $[0,1)$ converging to 1 we have that the  sequence $(\gamma(t_n))$ is restricted in the sense of Definition \ref{restricted}.
			Clearly this is equivalent to saying that $\gamma([0,1))$ is contained in a $K$-region with vertex $\xi$ and that $k_D(\gamma(t),\varphi(\D))\stackrel{t\to 1^-}\longrightarrow 0,$ where $\varphi$ is a complex geodesic with endpoint $\xi$.
			We say that  curve  $\gamma:[0,1)\to D$ with endpoint $\xi$ is {\sl strongly restricted} if $\gamma([0,1))$  is contained in a $K$-region with vertex $\xi$ and if there exists a (continuous) curve $\tilde \gamma\colon [0,1)\to \D$ such that
			\begin{equation}\label{stronglyrestricted}\lim_{t\to1^-}k_D(\gamma(t),\varphi(\tilde \gamma(t)))=0,
			\end{equation}
			where $\varphi$ is a complex geodesic with endpoint $\xi$.
		\end{definition}

		\begin{theorem}[Lindel\"{o}f Principle]\label{localLindelof}
			Let $D\subset\C^d$ be a $\C$-proper  convex domain and let $\xi\in\partial D$ be a point of locally finite type. Let $f\colon D\to \C$  be a holomorphic function which is $K$-bounded at $\xi$, and assume that there exists a strongly restricted curve $\gamma\colon[0,1)\to D$  such that $f(\gamma(t))\to L$ as $t\to 1$. Then $f$ has restricted $K$-limit $L$ at $\xi$.
		\end{theorem}

		\begin{lemma}\label{goodprodev}
			Let $D\subset\C^d$ be a $\C$-proper convex domain and let $\xi\in\partial D$ be a point of locally finite type.
			Let $(z_n)$ be a restricted sequence converging to $\xi$, and let $\varphi\colon \D\to D$ be a complex geodesic with endpoint $\xi$. Let $R>0$ be such that  $z_n\in A(\varphi,R)$ for all $n\geq 0$.
			Let  $(\zeta_n)$ be a sequence in $\D$ such that 
			$k_D(z_n,\varphi(\zeta_n))\to 0.$ Then for all $R'>R$ the sequence $(\varphi(\zeta_n))$ is eventually contained in $A(\varphi,R'$) and we have
			$$\lim_{n\to+\infty}k_{A(\varphi,R')}(z_n,\varphi(\zeta_n))=0.$$
		\end{lemma}

		\begin{remark}It follows from the previous lemma that, if $\gamma\colon[0,1)\to D$ is a strongly restricted curve such that $\gamma(t)\in A(\varphi,R)$ for all $t\in[0,1)$ and $\tilde{\gamma}\colon[0,1)\to\D$  is a curve satisfying \eqref{stronglyrestricted}, then for all $R'>R$ we have that 
			$\varphi(\tilde\gamma(t))$ is eventually contained in $A(\varphi,R')$ and
			$$\lim_{t\to 1^-}k_{A(\varphi,R')}(\gamma(t),\varphi(\tilde\gamma(t)))=0.$$
		\end{remark}

		\begin{proof}[Proof of Theorem \ref{localLindelof}]
			Let  $\varphi\colon \D\to D$ be a complex geodesic with endpoint $\xi$ and 
			let $\tilde \gamma\colon[0,1)\to \D$ be a curve satisfying \eqref{stronglyrestricted}.
			Let $R>0$ such that $\gamma(t)\in A(\varphi,R)$ for all $t\in[0,1)$, and 	let $R'>R$. 
			Since $f$ is $K$-bounded at $\xi$ there exists $S>0$ such that $f(A(\varphi,R'))\subset\subset S\D$, and thus by Lemma \ref{goodprodev}
			$$k_{S\D}(f(\gamma(t)),f(\varphi(\tilde\gamma(t))))\leq k_{A(\varphi,R')}(\gamma(t),\varphi(\tilde \gamma(t)))\stackrel{t\to 1^-}\longrightarrow 0,$$ 
			thus
			$$
			\lim_{t\to1^-}|f(\gamma(t))-f(\varphi(\tilde \gamma(t)))|=0.
			$$
			Since $f(\gamma(t))\to L$ as $t\to 1^-$ it follows that $f(\varphi(\tilde \gamma(t)))\to L$. Applying the classical Lindel\"{o}f Theorem to the holomorphic function $f\circ\varphi\colon \D\to \C$ we have that $f\circ\varphi$ has non-tangential limit $L$ at $1$.
			
			
			
			Let now $(z_n)$ be a  restricted sequence in $D$ converging to $\xi$, and let  $(\zeta_n)$ be a sequence in $\D$ such that 
			$k_D(z_n,\varphi(\zeta_n))\to 0.$
			Arguing as above we obtain that 
			\begin{equation}\label{LindEQ}
				\lim_{n\to+\infty}|f(z_n)-f(\varphi(\zeta_n))|=0.
			\end{equation}
			Since  $(\zeta_n)$ converges to $1$ non-tangentially, we have  $f(\varphi(\zeta_n))\to L$. By \eqref{LindEQ} we obtain $f(z_n)\to L$ as $n\to+\infty$, and we are done.
		\end{proof}

\end{document}
